\documentclass[11pt]{article}
\usepackage[T1]{fontenc}
\usepackage[utf8]{inputenc}
\usepackage{lmodern}
\usepackage[a4paper,margin=1in]{geometry}
\usepackage{amsmath,amssymb,amsthm,mathtools}
\usepackage{microtype,enumitem,booktabs}
\usepackage[colorlinks=true,linkcolor=blue,citecolor=blue,urlcolor=blue]{hyperref}
\newtheorem{theorem}{Theorem}[section]
\newtheorem{proposition}[theorem]{Proposition}
\newtheorem{lemma}[theorem]{Lemma}
\newtheorem{corollary}[theorem]{Corollary}
\theoremstyle{definition}
\newtheorem{assumption}[theorem]{Input}
\theoremstyle{remark}
\newtheorem{remark}[theorem]{Remark}
\newcommand{\Q}{\mathbb Q}
\newcommand{\Z}{\mathbb Z}
\newcommand{\F}{\mathbb F}
\newcommand{\C}{\mathbb C}
\DeclareMathOperator{\rank}{rank}
\DeclareMathOperator{\Res}{Res}
\DeclareMathOperator{\Gal}{Gal}
\DeclareMathOperator{\ord}{ord}
\DeclareMathOperator{\supp}{supp}
\title{Conditional Corollaries from the October 2026\newline OpenAI Mathematics Release}
\author{John Gregory George\\
\small Center for Economic Analysis, Middle Georgia State University}
\date{Version 1.0; prepared 7 October 2026}
\hypersetup{pdftitle={Conditional Corollaries from the October 2026 OpenAI Mathematics Release},pdfauthor={John Gregory George},pdfsubject={Conditional mathematical synthesis; AI contributions disclosed}}
\begin{document}
\maketitle
\begin{center}\small
Project conception and direction: John Gregory George.\\
Mathematical synthesis, proof drafting and source checks: ChatGPT (OpenAI).\\
Motivation and inspiration: Jarvis, George's local AI research system.
\end{center}
\begin{abstract}
We record selected conditional consequences of manuscripts in the October 2026 OpenAI mathematics release. A primitive-root prime lower bound and a quantitative progression-free-set estimate imply disjoint arithmetic-progression packings covering a proportion tending to one of the relevant primes in every sufficiently large dyadic interval. A related matching argument gives prime-polynomial differences. Low-Selmer-corank BSD and quadratic-twist rank statistics imply density-one full BSD over each fixed multiquadratic field; in degree four, at least one third in lower density of rational signed-squarefree parameters give full-BSD twists of rank at least two. Further consequences concern nearly orthogonal binary circulants, Hodge classes on products of K3-associated cubic fourfolds, and idempotents in a 2-adic completion. We state dependencies explicitly and prove the transfer arguments. The cited preprint theorems are assumed, not reproved or certified. Historical priority is not established.
\end{abstract}

\section{Status, scope and attribution}
This note records mathematical implications, not independent proofs of the source manuscripts. Every result below is conditional on the precise inputs cited in its statement or section. The major upstream theorems belong to their cited authors. The elementary packing, finite-intersection, decomposition and lifting mechanisms used here are standard; the purpose is to make the particular consequences and their dependencies explicit.

All OpenAI manuscript references are pinned to repository revision
\begin{center}\small\texttt{adc7f1241b42e322a6451854ab7e4b4c146bf78a}.\end{center}
The release describes mixed verification status. No Lean verification was rerun for this note. In particular, the official documentation for family 159 states that its selected formalized result is the divergent-reciprocal-sum consequence, excluding the quantitative progression-free-set estimate used in Section~\ref{sec:packing}; see \cite{formal159}. A formalization badge must not be interpreted as certification of that quantitative input.

The arguments were assembled during an AI-assisted review on 7 October 2026. A private preparation date is not a public disclosure timestamp. No assertion of first discovery is made: targeted searches and the absence of an explicitly stated corollary in inspected manuscripts do not establish exhaustive priority. This version has not undergone external peer review.

\paragraph{Contribution disclosure.}
George conceived and directed the cross-paper synthesis project and requested public documentation and attribution. ChatGPT (OpenAI) performed the source screening, proposed and developed the displayed conditional deductions, conducted mathematical interface and bibliographic checks, and drafted this note, using parallel AI review passes. Those passes are not independent human peer review. Jarvis is acknowledged for motivation and inspiration; no proof in this note is attributed to Jarvis. The human author is the responsible scholarly contact. AI contributions are disclosed explicitly rather than represented as human coauthorship. The institutional affiliation identifies the author; no institutional endorsement is implied.

\section{Packing arithmetic progressions of primitive-root primes}\label{sec:packing}
For a fixed integer $a$ that is neither $-1$ nor a square, put
\[
\mathcal A_a(N)=\{p\text{ prime}:N<p<2N,\ p\nmid a,\ \ord_p(a)=p-1\}.
\]
Let $r_k(N)$ be the maximum cardinality of a subset of $[N]=\{1,\ldots,N\}$ containing no $k$-term arithmetic progression with positive common difference.

\begin{assumption}[Primitive-root abundance]\label{in:artin}
For each fixed admissible $a$, there is $c_a>0$ such that
\[|\mathcal A_a(N)|\ge c_aN/(\log N)^2\]
for all sufficiently large $N$. This is the portion of \cite[Theorem 1.1, p.~2]{artin} used here.
\end{assumption}
\begin{assumption}[Quantitative progression-free bound]\label{in:ap}
For each fixed integer $k\ge3$, there exist positive constants $C_k,c_k,\varepsilon_k$ such that
\[r_k(N)\le C_kN\exp\{-c_k(\log N)^{\varepsilon_k}\}.\]
We use \cite[Theorem 1.1, p.~4]{ap}.
\end{assumption}

\begin{lemma}[Maximal-packing remainder]\label{lem:packing}
Every $A\subseteq[N]$ admits a family of pairwise vertex-disjoint $k$-term arithmetic progressions leaving at most $r_k(N)$ elements uncovered.
\end{lemma}
\begin{proof}
Choose a maximal such family, or repeatedly delete a progression while one exists. The uncovered set contains no progression, since otherwise the family could be enlarged. Its size is therefore at most $r_k(N)$.
\end{proof}

\begin{theorem}[Near-complete progression packing]\label{thm:artin}
Assume Inputs~\ref{in:artin} and~\ref{in:ap}. For each fixed admissible $a$ and fixed $k\ge3$, every sufficiently large integer $N$ admits disjoint $k$-term arithmetic progressions in $\mathcal A_a(N)$ whose uncovered set $R_N$ satisfies
\[
|R_N|\le C_kN e^{-c_k(\log N)^{\varepsilon_k}},\qquad
\frac{|R_N|}{|\mathcal A_a(N)|}\longrightarrow0.
\]
In particular, eventually there are at least $c_aN/(2k(\log N)^2)$ disjoint progressions.
\end{theorem}
\begin{proof}
Translation by $-N$ puts $\mathcal A_a(N)$ inside $[N]$ and preserves progressions. Apply Lemma~\ref{lem:packing}, then divide by Input~\ref{in:artin}:
\[
\frac{|R_N|}{|\mathcal A_a(N)|}
\le\frac{C_k}{c_a}(\log N)^2e^{-c_k(\log N)^{\varepsilon_k}}=o(1).
\]
Eventually at least half the vertices are covered; division by $k$ gives the final bound.
\end{proof}

\begin{remark}[Scope and prior work]
The progression length is fixed before $N\to\infty$, and different progressions may have different common differences. This is an existence result without a practical construction or useful threshold claim. Input~\ref{in:artin} alone does not give positive relative density among primes or divergence of the reciprocal sum; neither the usual relative-density Green--Tao theorem nor the selected formalized reciprocal-sum statement can replace Input~\ref{in:ap} without further work. Frei and Tinkov\'a \cite[Theorem 1.2]{constellations} establish stronger primitive-root constellation asymptotics under Hooley-type Riemann hypotheses. The present narrower implication uses the newly stated Inputs~\ref{in:artin}--\ref{in:ap}; it does not claim that the packing method itself is new.
\end{remark}

\subsection{Prime-polynomial difference matchings}
\begin{proposition}\label{prop:matching}
Assume Input~\ref{in:artin} and \cite[Theorem 1.1, p.~1]{polydiff}. Fix an integer $j\ge2$. In every sufficiently large interval $(N,2N)$, disjoint pairs of primes in $\mathcal A_a(N)$, each having difference $(q-1)^j$ for some prime $q$, cover a proportion tending to one of $\mathcal A_a(N)$.
\end{proposition}
\begin{proof}
The polynomial $h(t)=(t-1)^j$ is prime-intersective: the unit $1$ is a root modulo every positive integer. The cited theorem bounds the size of any subset of $[N]$ without a nonzero difference $h(q)$, $q$ prime, by $C_hN^{1-c_h}$ for constants $C_h,c_h>0$. Form the graph on $\mathcal A_a(N)$ joining primes with such a difference, and take a maximal matching. Unmatched vertices form an independent set. After translation its size is at most $C_hN^{1-c_h}$, which divided by Input~\ref{in:artin} is $O_{a,j}((\log N)^2N^{-c_h})=o(1)$.
\end{proof}
The witnesses $q$ need not be distinct. Proposition~\ref{prop:matching} and Theorem~\ref{thm:artin} do not imply a progression with common difference $(q-1)^j$: their witnesses need not coincide.

\section{BSD for rational twists over a fixed multiquadratic field}
For an elliptic curve $E/\Q$, write $E^{(d)}$ for its quadratic twist and set
\[
\mathcal D(X)=\{d\in\Z:0<|d|\le X,\ d\text{ squarefree}\}.
\]
Every density and mean below uses this signed rational parameter set. Write $r_F(d)=\rank F^{(d)}(\Q)$ and $a_F(d)=\ord_{s=1}L(F^{(d)},s)$.

We assume the following portions of the cited release:
\begin{enumerate}[label=(\roman*)]
\item For each fixed $F/\Q$, there is a density-one set on which $r_F(d)=a_F(d)\in\{0,1\}$ and the whole Tate--Shafarevich group is finite; each rank occurs with density $1/2$ \cite[Theorems 1.1--1.2, p.~2]{density}.
\item Full $q$-power Selmer corank zero or one implies the full BSD formula for elliptic curves over $\Q$ \cite[Theorem 1.1, p.~3]{bsd}.
\item For the separate all-parameter mean assertion only, the analytic and algebraic mean ranks of $F^{(d)}$ tend to $1/2$ \cite[Theorem 1.1 and Corollary 1.3, pp.~1--2]{mean}.
\end{enumerate}
The compatibility of full BSD with products, isogenies and Weil restriction is standard; a convenient primary reference is \cite[Proposition 6]{vanbommel}.

\begin{theorem}[Fixed multiquadratic base change]\label{thm:bsd}
Fix $E/\Q$ and a multiquadratic field $K/\Q$ of degree $m=2^r$. Assuming (i)--(ii), there is a density-one set $G\subseteq\mathcal D(\infty)$ for which $E^{(d)}/K$ satisfies the full BSD formula and has rank at most $m$. Assuming also (iii), both its mean algebraic rank and its mean analytic rank over all $d\in\mathcal D(X)$ tend to $m/2$.
\end{theorem}
\begin{proof}
For each character $\chi$ of $\Gal(K/\Q)$, let $F_\chi=E\otimes\chi$, a fixed elliptic curve over $\Q$. Apply (i) to each of these $m$ curves, using the same parameter $d$, and intersect the resulting good sets. Since $m$ is fixed, their intersection $G$ has density one. Finite Tate--Shafarevich group identifies the $2$-power Selmer corank with the Mordell--Weil rank, so (ii) gives full BSD for every $F_\chi^{(d)}$ on $G$.

The character decomposition gives a $\Q$-isogeny
\[
\Res_{K/\Q}(E^{(d)}_K)\sim_{\Q}\prod_{\chi}F_\chi^{(d)}.
\]
The full BSD compatibilities therefore give full BSD for $E^{(d)}/K$ on $G$. They transport all factors of the leading-term formula, including regulators and isogeny indices, not just the $L$-values. Moreover, for every $d$,
\begin{align}
R(d):=\rank E^{(d)}(K)&=\sum_\chi r_{F_\chi}(d),\label{eq:rank}\\
\ord_{s=1}L(E^{(d)}/K,s)&=\sum_\chi a_{F_\chi}(d),\label{eq:analytic}
\end{align}
where the second identity follows from the $L$-function factorization. On $G$ the first sum is at most $m$. Under (iii), taking means in the two exact identities proves the all-parameter limits by a finite summation.
\end{proof}

\begin{corollary}[Positive lower density of higher-rank full-BSD twists]\label{cor:highrank}
Assume only (i)--(ii) and fix $E,K$ as above with $m\ge2$. Then
\[
\liminf_{X\to\infty}\frac{\#\{d\in\mathcal D(X): E^{(d)}/K\text{ satisfies full BSD},\ R(d)\ge m/2\}}{\#\mathcal D(X)}\ge\frac{2}{m+2}.
\]
In particular, for any fixed biquadratic $K/\Q$, at least one third in lower density of these rational parameters give full-BSD twists of rank at least two over $K$.
\end{corollary}
\begin{proof}
For each $\chi$, removing the zero-density complement of $G$ does not change the half-density of the event $r_{F_\chi}(d)=1$. Since each component is zero or one on $G$, \eqref{eq:rank} yields
\[
\frac1{\#\mathcal D(X)}\sum_{d\in G\cap\mathcal D(X)}R(d)
=\sum_\chi\frac{\#\{d\in G\cap\mathcal D(X):r_{F_\chi}(d)=1\}}{\#\mathcal D(X)}\longrightarrow\frac m2.
\]
Set $b=m/2$ and $p_X=\#\{d\in G\cap\mathcal D(X):R(d)\ge b\}/\#\mathcal D(X)$. On the rest of $G$ the integer $R(d)$ is at most $b-1$, and on $G$ it is always at most $m$. Thus
\[
\frac1{\#\mathcal D(X)}\sum_{d\in G\cap\mathcal D(X)}R(d)
\le (b-1)+(b+1)p_X.
\]
Taking lower limits proves $\liminf p_X\ge1/(b+1)=2/(m+2)$. Every parameter counted by $p_X$ satisfies full BSD by Theorem~\ref{thm:bsd}.
\end{proof}

\begin{remark}
No independence of the component ranks is used. Assumption (iii) is unnecessary for Corollary~\ref{cor:highrank}; it is needed only for the separate mean over all parameters. Density one alone would not justify that unrestricted mean. The conclusions concern rational twists of a curve descending from $\Q$, with $K$ fixed, rather than all elliptic curves over $K$, all intrinsic $K$-twists, or distinct $K$-isomorphism classes. No uniformity in growing $K$ or exact rank distribution is asserted.
\end{remark}

\section{Nearly orthogonal binary circulants}
Assume \cite[Theorem 1, p.~1]{ultraflat}: for every sufficiently large integer $N$ there are signs $\epsilon_j\in\{-1,1\}$ such that
\[
P_N(z)=\sum_{j=0}^{N-1}\epsilon_jz^j,\qquad
\sup_{|z|=1}\left|\frac{|P_N(z)|}{\sqrt N}-1\right|=\eta_N\longrightarrow0.
\]
\begin{proposition}
There exist real circulant $\{-1,1\}$ matrices $C_N$ in every sufficiently large order such that
\[
\kappa_2(C_N)\longrightarrow1,\qquad
\left\|\frac{C_NC_N^*}{N}-I\right\|_{\mathrm{op}}\longrightarrow0.
\]
If \cite[Theorem 1.1, p.~1]{circulant} is also assumed, exact equality $C_NC_N^*=NI$ is impossible for $N>4$.
\end{proposition}
\begin{proof}
Take the first row of $C_N$ to be the coefficients of $P_N$. Fourier diagonalization of circulants shows that its singular values are $|P_N(\omega_N^j)|$, where $\omega_N=e^{2\pi i/N}$, up to an immaterial conjugation convention. They lie between $(1-\eta_N)\sqrt N$ and $(1+\eta_N)\sqrt N$. Consequently
\[
\kappa_2(C_N)\le\frac{1+\eta_N}{1-\eta_N},\qquad
\left\|C_NC_N^*/N-I\right\|_{\mathrm{op}}\le2\eta_N+\eta_N^2.
\]
The exact nonexistence statement is the cited circulant-Hadamard theorem.
\end{proof}
The flat-polynomial/circulant method is already explicit in Steinerberger \cite{steinerberger}. The strengthened asymptotic follows from the newly claimed ultraflat input; neither a new method nor an efficient construction of the signs is claimed.

\section{Products of cubic fourfolds with geometric K3 categories}
\label{sec:cubic-products}

The assertions in this section are conditional on the specified theorem in the
OpenAI release.  The argument supplies a corollary of that theorem through an
established motivic comparison; it does not independently verify the released
proof or establish priority for the corollary.

\begin{assumption}[K3-product input]
\label{ass:k3-products}
Theorem~1.1 of \cite{k3} holds: for every finite list of smooth projective
complex K3 surfaces $S_1,\ldots,S_m$, allowing repeated factors, and every
integer $p$, the rational cycle-class map
\[
 \mathrm{CH}^{p}(S_1\times\cdots\times S_m)_{\mathbb Q}
 \longrightarrow
 H^{2p}(S_1\times\cdots\times S_m,\mathbb Q)
       \cap H^{p,p}(S_1\times\cdots\times S_m,\mathbb C)
\]
is surjective.  The intersection is taken after extension of scalars to
complex cohomology.
\end{assumption}

\begin{proposition}[Rational Hodge for products of geometric-K3 cubics]
\label{prop:cubic-products}
Assume Assumption~\ref{ass:k3-products}.  Let $X_1,\ldots,X_r$ be smooth
complex cubic fourfolds, with $r\geq 1$.  For each $i$, suppose there is a
smooth projective complex K3 surface $S_i$ and a Fourier--Mukai equivalence
\[
 \operatorname{Ku}(X_i)\simeq D^b(\operatorname{Coh}S_i),
\]
where $\operatorname{Ku}(X_i)$ is the Kuznetsov component and the K3 category
on the right is untwisted.  Then the rational Hodge conjecture holds in every
codimension on $X_1\times\cdots\times X_r$.  In particular, it holds on every
self-power of each such cubic fourfold.
\end{proposition}

\begin{proof}
We work with Chow motives over $\mathbb C$ with rational coefficients, using
the convention that $\mathbb Q(-1)$ has its Betti realization in degree two.
The algebraic and transcendental Chow--K\"unneth decomposition for a smooth
cubic fourfold, as recalled in \cite[\S5.2]{fuvial}, gives
\[
 h(X_i)\simeq
 \mathbb Q(0)\oplus\mathbb Q(-1)
 \oplus\mathbb Q(-2)^{\oplus a_i}
 \oplus h^4_{\mathrm{tr}}(X_i)
 \oplus\mathbb Q(-3)\oplus\mathbb Q(-4),
\]
where $a_i$ is the dimension of its rational algebraic middle cohomology.
Fu--Vial's theorem \cite[Theorem~3, equivalently Theorem~7.2]{fuvial}
applies to the stated Fourier--Mukai equivalence and gives an isomorphism
of rational Chow motives
\[
 h^4_{\mathrm{tr}}(X_i)(2)
       \simeq h^2_{\mathrm{tr}}(S_i)(1).
\]
Thus the non-Tate summand of $h(X_i)$ is isomorphic to
$h^2_{\mathrm{tr}}(S_i)(-1)$, itself an algebraic direct summand of a Tate
twist of $h(S_i)$.  These isomorphisms and their inverses are algebraic
correspondences, not merely isomorphisms of Hodge structures.

Tensor the displayed decompositions for all $i$.  Each summand of
$h(X_1\times\cdots\times X_r)$ is isomorphic to an algebraic direct
summand of a Tate twist of $h(\prod_{i\in I}S_i)$ for some subset
$I\subseteq\{1,\ldots,r\}$; for $I=\varnothing$ it is a Tate motive.
Assumption~\ref{ass:k3-products} applies to every nonempty such product.
Its conclusion passes to algebraic direct summands: if $\pi$ is the
corresponding algebraic projector and $\alpha$ is a Hodge class in its
image, first represent $\alpha$ by an algebraic cycle $z$ on the full
product and then use $\pi_*z$.  Tate twists only shift the degree and
codimension.  Finally transport these cycle classes through the inverse
algebraic correspondences and sum the components.  This represents every
rational Hodge class on the cubic product by an algebraic cycle.
\end{proof}

\begin{proposition}[Application to the released irrational cubics]
\label{prop:irrational-cubics}
Assume Assumption~\ref{ass:k3-products} and Theorem~1.1 of \cite{cubic}.
For every sufficiently large admissible Hassett discriminant $d$, a very
general smooth cubic fourfold $X$ in the Hassett divisor $\mathcal C_d$
is irrational and satisfies the rational Hodge conjecture on every
self-power.  More generally, the rational Hodge conjecture holds on every
finite product of cubics chosen in these very-general loci, even when the
discriminants differ.
\end{proposition}

\begin{proof}
The cited cubic theorem supplies irrationality and an exact complex-linear
equivalence of $\operatorname{Ku}(X)$ with the ordinary derived category
of a smooth projective K3 surface.  For this untwisted target the equivalence
is Fourier--Mukai, as recorded following Theorem~3 of \cite{fuvial}.
Proposition~\ref{prop:cubic-products} therefore applies to any finite list
of these cubics.
\end{proof}

\begin{remark}
Here admissible means $d>6$, $d\equiv 0$ or $2\pmod 6$, with $4\nmid d$,
$9\nmid d$, and no odd prime congruent to $2\pmod 3$ dividing $d$, as in
\cite{cubic}.  The discriminant threshold is ineffective, and ``very
general'' is interpreted separately in each $\mathcal C_d$.  The proposition
does not assert the Hodge conjecture for all cubic fourfolds, the full Hodge
conjecture, or stable irrationality.  The correspondence bridge is an
established theorem of Fu and Vial; the release supplies the K3-product and
irrationality inputs.  Novelty of the combined corollary is not established.
\end{remark}

\section{Idempotent lifting with necessarily infinite support}
\label{sec:adic-support}

Write $\mathbb Z_2$ for the ring of $2$-adic integers and
$\mathbb F_2$ for the field with two elements.  For a discrete group $G$,
$R[G]$ denotes its ordinary group ring: all its elements have finite group
support.  The completion below completes coefficients; it does not replace
$G$ by a profinite group.

\begin{assumption}[The two group-ring inputs]
\label{ass:group-rings}
The following statements from the release hold.
\begin{enumerate}
\item Theorem~1.1 of \cite{group2}: there exist a finitely presented
      torsion-free group $G$ with a finite two-dimensional classifying
      complex and $a,b,c\in\mathbb F_2[G]$ satisfying
      $ab=1$, $ac=0$, and $c\neq 0$.
\item Corollary~1.2(ii) of \cite{idempotent}: for every torsion-free
      discrete group $H$ and every commutative unital characteristic-zero
      domain $R$, the only idempotents of $R[H]$ are $0$ and $1$.
\end{enumerate}
The second statement concerns scalar ring elements, not arbitrary matrices.
\end{assumption}

\begin{proposition}[Completion creates an idempotent requiring infinite support]
\label{prop:adic-support}
Assume Assumption~\ref{ass:group-rings}, and take its group $G$ and elements
$a,b,c$.  Set
\[
 A=\mathbb Z_2[G],\qquad
 \widehat A=\varprojlim_{n\geq1}(\mathbb Z/2^n\mathbb Z)[G],
 \qquad e=1-ba\in\mathbb F_2[G].
\]
Then $e$ is a nontrivial idempotent and lifts to an idempotent of
$\widehat A$.  The ring $A$ has no nontrivial idempotents, and every
idempotent of $\widehat A$ reducing to $e$ has infinite support in $G$.
For every compatible sequence of idempotents
\[
 E_n\in(\mathbb Z/2^n\mathbb Z)[G],\qquad E_1=e,
\]
the finite numbers $|\operatorname{supp}E_n|$ tend to infinity as
$n\to\infty$.
\end{proposition}

\begin{proof}
First $ba\neq 1$, since otherwise
$c=(ba)c=b(ac)=0$.  As $ab=1$, we have $(ba)^2=b(ab)a=ba$, and hence
$e^2=e$.  The element $e$ is nonzero.  It is not $1$: if $ba=0$, then
$b=b(ab)=(ba)b=0$, contradicting $ab=1$.

Choose a finite-support integral lift $x_0\in\mathbb Z[G]\subset A$ of
$e$.  Define recursively
\[
 x_{j+1}=3x_j^2-2x_j^3,\qquad q_j=x_j^2-x_j.
\]
The following identities hold in any unital ring:
\[
 q_{j+1}=q_j^2(4q_j-3),\qquad
 x_{j+1}-x_j=(1-2x_j)q_j.
\]
Indeed, they are polynomial identities in the single element $x_j$ and
central integer scalars, so noncommutativity of the ambient ring causes
no issue.  Since $q_0\in2A$, induction gives
$q_j\in2^{2^j}A$.  The second identity makes $(x_j)$ a $2$-adic Cauchy
sequence.  Its limit $E\in\widehat A$ satisfies $E^2=E$ and reduces to
$e$ modulo $2$, so it is nontrivial.

For clarity, an element of $\widehat A$ can be identified with a
coefficient family $(u_g)_{g\in G}$ in $\mathbb Z_2$ such that, for every
$n$, only finitely many $u_g$ are nonzero modulo $2^n$.  Its support is
$\{g:u_g\neq0\}$.  Such an element has finite support exactly when it
lies in the natural embedded copy of $A$.

The second input, with $R=\mathbb Z_2$ and $H=G$, excludes every
nontrivial idempotent of $A$.  Thus any idempotent lift of $e$ in
$\widehat A$ has infinite support.  For a compatible sequence $(E_n)$,
its limit is such a lift.  Compatibility makes the sets
$\operatorname{supp}E_n$ nested, and their union is the infinite support
of the limit.  A nested sequence of finite sets with infinite union has
cardinalities tending to infinity, proving the last assertion.
\end{proof}

\begin{proposition}[Failure of direct finiteness after completion]
\label{prop:adic-direct-finiteness}
The first input of Assumption~\ref{ass:group-rings} alone implies that
$\widehat A$ is not directly finite: there exist $U,V\in\widehat A$
with $UV=1$ and $VU\neq1$.
\end{proposition}

\begin{proof}
Choose finite-support lifts $A_0,B_0\in A$ of $a,b$.  Then
$A_0B_0=1+2C$ for some $C\in A$.  The geometric series
$\sum_{j\geq0}(-2C)^j$ converges in $\widehat A$ and is a two-sided
inverse of $A_0B_0$.  Put
$U=A_0$ and $V=B_0(A_0B_0)^{-1}$.  Then $UV=1$, whereas
$VU$ reduces to $ba\neq1$ modulo $2$.
\end{proof}

\begin{remark}
The obstruction is finite support, not the existence of lifts at finite
precision.  No quantitative lower bound on support growth is asserted.
The lifting polynomials and the geometric-series argument are standard
completion methods; the conditional separation combines them with the
specified release inputs.  Its prior publication status is not established.
\end{remark}


\section{Interpretation and priority boundary}
These propositions are short transfer arguments from powerful upstream results. They do not establish the Riemann hypothesis, unrestricted BSD, the unrestricted Hodge conjecture, or an additional Millennium Prize solution. Acceptance of an implication does not certify its premises. Some consequences may already be known, or may be considered immediate by specialists; this note makes no exhaustive novelty claim.

The intended public record identifies exactly what was written, by whom and with which disclosed assistance, at the time a hosting service actually publishes it. It cannot certify absence of earlier work. Corrections, stronger source theorems, and overlooked precedents should be recorded in dated revisions with the earlier versions retained.

\begingroup\small
\begin{thebibliography}{99}

\bibitem{artin} OpenAI. \emph{Primitive roots for every admissible integer base}. Preprint, 2026. Release identifier 029-01. \href{https://github.com/openai/math/blob/adc7f1241b42e322a6451854ab7e4b4c146bf78a/preprints/Primitive-roots-for-every-admissible-integer-base-October-4-2026/primitive-roots-all-integer-bases.pdf}{Pinned manuscript}.

\bibitem{ap} OpenAI. \emph{Quasipolynomial Bounds for Arithmetic Progressions}. Preprint, 2026. Release identifier 159-01. \href{https://github.com/openai/math/blob/adc7f1241b42e322a6451854ab7e4b4c146bf78a/preprints/Quasipolynomial-Bounds-for-Arithmetic-Progressions-September-23-2026/paper.pdf}{Pinned manuscript}.

\bibitem{polydiff} OpenAI. \emph{A Power Saving for Polynomial Differences at Prime Arguments}. Preprint, 2026. Release identifier 182-02. \href{https://github.com/openai/math/blob/adc7f1241b42e322a6451854ab7e4b4c146bf78a/preprints/A-Power-Saving-for-Polynomial-Differences-at-Prime-Arguments-October-5-2026/prime-argument-polynomial-differences.pdf}{Pinned manuscript}.

\bibitem{bsd} OpenAI. \emph{Exact Birch–Swinnerton-Dyer Formula from Low Selmer Corank}. Preprint, 2026. Release identifier 002-01. \href{https://github.com/openai/math/blob/adc7f1241b42e322a6451854ab7e4b4c146bf78a/preprints/Exact-Birch-Swinnerton-Dyer-Formula-from-Low-Selmer-Corank-October-3-2026/exact-bsd-low-selmer-corank.pdf}{Pinned manuscript}.

\bibitem{density} OpenAI. \emph{Goldfeld's analytic density conjecture and the 2-converse for elliptic curves}. Preprint, 2026. Release identifier 006-01. \href{https://github.com/openai/math/blob/adc7f1241b42e322a6451854ab7e4b4c146bf78a/preprints/Goldfelds-analytic-density-conjecture-and-the-2-converse-for-elliptic-curves-September-23-2026/paper.pdf}{Pinned manuscript}.

\bibitem{mean} OpenAI. \emph{The mean analytic rank of quadratic twists of elliptic curves}. Preprint, 2026. Release identifier 006-02. \href{https://github.com/openai/math/blob/adc7f1241b42e322a6451854ab7e4b4c146bf78a/preprints/The-mean-analytic-rank-of-quadratic-twists-of-elliptic-curves-September-23-2026/paper.pdf}{Pinned manuscript}.

\bibitem{ultraflat} OpenAI. \emph{Ultraflat real Littlewood polynomials}. Preprint, 2026. Release identifier 076-01. \href{https://github.com/openai/math/blob/adc7f1241b42e322a6451854ab7e4b4c146bf78a/preprints/Ultraflat-real-Littlewood-polynomials-October-5-2026/ultraflat-real-littlewood-polynomials.pdf}{Pinned manuscript}.

\bibitem{circulant} OpenAI. \emph{The circulant Hadamard conjecture}. Preprint, 2026. Release identifier 179-01. \href{https://github.com/openai/math/blob/adc7f1241b42e322a6451854ab7e4b4c146bf78a/preprints/The-circulant-Hadamard-conjecture-September-23-2026/paper.pdf}{Pinned manuscript}.

\bibitem{k3} OpenAI. \emph{The rational Hodge conjecture for products of K3 surfaces}. Preprint, 2026. Release identifier 032-01. \href{https://github.com/openai/math/blob/adc7f1241b42e322a6451854ab7e4b4c146bf78a/preprints/The-rational-Hodge-conjecture-for-products-of-K3-surfaces-October-4-2026/hodge-conjecture-products-k3.pdf}{Pinned manuscript}.

\bibitem{cubic} OpenAI. \emph{Irrational cubic fourfolds with geometric K3 categories}. Preprint, 2026. Release identifier 054-01. \href{https://github.com/openai/math/blob/adc7f1241b42e322a6451854ab7e4b4c146bf78a/preprints/Irrational-cubic-fourfolds-with-geometric-K3-categories-September-24-2026/Irrational-cubic-fourfolds-with-geometric-K3-categories-September-24-2026.pdf}{Pinned manuscript}.

\bibitem{group2} OpenAI. \emph{A Torsion-Free Group Algebra That Is Not Directly Finite}. Preprint, 2026. Release identifier 197-01. \href{https://github.com/openai/math/blob/adc7f1241b42e322a6451854ab7e4b4c146bf78a/preprints/A-Torsion-Free-Group-Algebra-That-Is-Not-Directly-Finite-October-4-2026/direct-finiteness.pdf}{Pinned manuscript}.

\bibitem{idempotent} OpenAI. \emph{The Bass trace conjecture and the characteristic-zero Kaplansky idempotent conjecture}. Preprint, 2026. Release identifier 207-02. \href{https://github.com/openai/math/blob/adc7f1241b42e322a6451854ab7e4b4c146bf78a/preprints/The-Bass-trace-conjecture-for-complex-group-rings-September-24-2026/The-Bass-trace-conjecture-for-complex-group-rings-September-24-2026.pdf}{Pinned manuscript}.

\bibitem{constellations} Christopher Frei and Magdal\'ena Tinkov\'a. \emph{Linear constellations in primes with arithmetic restrictions}. With an appendix by Christopher Frei, Joachim K\"onig and Magdal\'ena Tinkov\'a. arXiv:2408.17441v3, 2025. \url{https://arxiv.org/abs/2408.17441v3}.

\bibitem{vanbommel} Raymond van Bommel. \emph{The Birch and Swinnerton-Dyer conjecture for an elliptic curve over $\mathbb Q(\sqrt[4]{5})$}. Archiv der Mathematik 114 (2020), 159--169. Proposition 6. \url{https://doi.org/10.1007/s00013-019-01383-w}.

\bibitem{fuvial} Lie Fu and Charles Vial. \emph{Cubic fourfolds, Kuznetsov components and Chow motives}. Author manuscript dated 5 September 2023, Theorem 3 (=Theorem 7.2). \url{https://www.math.uni-bielefeld.de/~vial/CubicFourfolds.pdf}.

\bibitem{steinerberger} Stefan Steinerberger. \emph{A Note on Approximate Hadamard Matrices}. arXiv:2402.13202, 2024. \url{https://arxiv.org/abs/2402.13202}.

\bibitem{formal159} OpenAI. \emph{Erd\H{o}s\textquotesingle s reciprocal-sum conjecture and quasipolynomial Szemer\'edi bounds: scope of the Lean formalization}. Repository documentation, 2026. \href{https://github.com/openai/math/blob/adc7f1241b42e322a6451854ab7e4b4c146bf78a/lean/docs/159.md}{Pinned scope statement}.

\end{thebibliography}

\endgroup
\end{document}
